\subsection{Retour longitudinal et détermination de l’échelle aréale}
\label{rev5:longitud}

L'unité dimensionnelle de l'aire admet une construction explicite à partir du même état enrichi. Nous écrirons \(L_\alpha\) pour la longueur physique obtenue ; l'entier \(L\) qui désigne la longueur d'un registre dans la section~\ref{sec:compat-lectores-formas} conserve son sens combinatoire. La distinction sépare le cardinal d'une représentation, l'étendue d'une arête et le rayon de sa réalisation.

La composition utilise les coordonnées régionales et le registre déjà
engendrés par APP--TRIT--TPK. APP conserve les deux feuillets avec leurs restes
et quotients ; TRIT oriente le régime local ; TPK prolonge l’état,
en conservant retenue, frontière, route et mémoire. La chaîne
\(w_6\to w_{12}\to w_{18}\to w_{24}\to w_{30}\to R_{36}\to G_9\)
et l’ordre \(U_t\to\Gamma_9\to K_{\rm ph}\) maintiennent le
retour de phase avec accroissement de mémoire. La construction conjointe
du continu conserve ses aspects spectral, solénoïdal, de jauge,
cohomologique et déterminantal. La longueur qui suit est une lecture
dimensionnelle ultérieure de cette structure et de ses sorties corrélées
\(\pi_{\rm HMT},\varphi_{\rm HMT},e_{\rm HMT},\alpha_{\rm HMT}\).

Fixons le registre marqué
\[
 \Omega=\{(j,k,q,u,v):j,k\in\mathbb Z/12\mathbb Z,\quad
 k-j\in\{0,3,6,9\},\quad1\le q\le8,\quad
 u<v,\quad u,v\in\{1,2,4,5,7,8\}\}.
\]
La soustraction est calculée dans ce groupe cyclique. Son cardinal est
\(N_\Omega=12\cdot4\cdot8\binom62=5760\).
Le facteur \(120\) provient de l’espace octadique des incidences
point--paire déjà construit ; les douze positions et les quatre
positions de chaque orbite conservent leurs marques respectives. Cette
résolution élargit le porteur du retour de Hadamard. Dans les espaces
hermitiens de départ et d’arrivée, de dimension \(N_\Omega\), soit
\(B\) l’inverseur transporté, avec
\(B^*B=BB^*=I/4\). L’identité locale
\(H_4^*H_4=4I\) produit cette normalisation pour l’inverseur ;
son extension orthogonale et le transport des marques la conservent.

Dans l'espace d'arrivée, soient \(u_\Omega=N_\Omega^{-1/2}\mathbf1\), \(P=u_\Omega u_\Omega^*\), \(C=(I-P)B\) et \(M=PB\). Le transport complet \(C+M=B\) est conservé. Les projecteurs atomiques \(P_i=e_ie_i^*\), déterminés par l'inscription, définissent l'espérance diagonale \(\Delta_\Omega(T)=\sum_iP_iTP_i\).

\begin{proposition}[Réponse enregistrée et complément conservé]
\label{rev5:retorno}
Parmi les applications linéaires vers l’algèbre diagonale qui fixent
les \(P_i\) et sont bimodulaires par rapport à cette algèbre,
\(\Delta_\Omega\) est unique. En outre,
\begin{equation}
 \Delta_\Omega(CC^*)=r_\Omega I,
 \qquad r_\Omega=\frac{N_\Omega-1}{4N_\Omega}
 =\frac{5759}{23040},
 \qquad \Delta_\Omega(MM^*)=\frac{I}{4N_\Omega}.
 \label{rev5:retorno-escalar}
\end{equation}
Les deux réponses ont pour somme \(I/4\).
\end{proposition}
\begin{proof}
Pour l’unité matricielle \(E_{ij}\), la bimodularité impose
\(E(E_{ij})=P_iE(E_{ij})P_j\). Comme l’image est diagonale,
ce terme vaut zéro si \(i\ne j\) ; si \(i=j\), fixer
\(P_i\) détermine son image. Les unités matricielles constituent
une base et déterminent l’application. D’autre part,
\(CC^*=(I-P)/4\), \(MM^*=P/4\), et chaque coefficient diagonal
de \(P\) est \(1/N_\Omega\). Appliquer l’espérance démontre
les trois identités.
\end{proof}

L'inscription isométrique \(We_i=e_i\otimes e_i\) réalise \(W^*(T\otimes I)W=\Delta_\Omega(T)\). L'inscription complète conserve l'information du registre ; l'espérance extrait sa réponse diagonale. Le complément \(M\) demeure dans le transport et possède la réponse spécifique de \eqref{rev5:retorno-escalar}. Cette séparation identifie l'observable qui intervient dans la longueur et la composante complémentaire conservée.

La forme normale de Smith de \(H_4\) est
\(\operatorname{diag}(1,2,2,4)\) ; son conoyau a pour ordre
seize. Le produit du scalaire engendré
\(\alpha=\alpha_{\rm HMT}\), indexé par les classes de ce
conoyau, définit le lecteur multiplicatif fini
\[
 n_H:=\prod_{\bar v\in\operatorname{coker}_{\mathbb Z}(H_4)}
 \alpha=\alpha^{16}.
\]
Soit \(L_*>0\) une section de référence
radiale de la droite dimensionnelle de longueur. La composition longitudinale
agit linéairement sur une cochaîne orientée \(\beta_*\) :
\begin{equation}
 \mathscr T_L\beta_*
 =\alpha^{16}(\Delta_\Omega(CC^*)\otimes I)\beta_*
 =q_\alpha\beta_*,\qquad
 q_\alpha=\alpha^{16}r_\Omega,\qquad
 L_\alpha=q_\alpha L_*.
 \label{rev5:longitud-generada}
\end{equation}
La formule fixe une composition concrète de lecteurs. La linéarité conserve le type de longueur et son orientation ; le coefficient provient du retour enregistré et de la norme locale. Par exemple, la réponse diagonale d'une source \(a e_i\) est \(a r_\Omega\), tandis que la norme \(\|C^*ae_i\|=|a|\sqrt{r_\Omega}\) constitue une autre lecture du même transport.

Si \(\beta_*(\bar a)=-U(a)\beta_*(a)\), multiplier par
\(q_\alpha\) conserve l’inversion orientée. Si une arête est
subdivisée en neuf tronçons compatibles,
\(\beta_*(a)=\sum_jU_j^{-1}\beta_*(a_j)\), la même
multiplication conserve cette identité et la concaténation de mémoire.
Dans un élargissement auxiliaire du registre, on transporte
\(P\), \(CC^*\) et \(\Delta_\Omega\) par tensorisation
avec l’identité. Ainsi, \(r_\Omega\) demeure ; le cardinal d’un
maillage raffiné et le cardinal du registre marqué ont des fonctions
distinctes.

\subsection{Réciprocité radiale et reconnaissance de la longueur de Planck}
\label{rev5:radial}

La lecture circulaire du calendrier de cent huit avancées conserve
\(108\ell_0=2\pi_{\rm HMT}L_\alpha\), avec
\(\ell_0=ct_0\). Par conséquent,
\begin{equation}
 t_0=\frac{\pi_{\rm HMT}L_\alpha}{54c},\qquad
 \omega=\frac{c}{L_\alpha},\qquad
 Q_{\rm clk}=\hbar_{\rm ret}\omega
 =\frac{\hbar_{\rm ret}c}{L_\alpha}.
 \label{rev5:reloj}
\end{equation}
L’action retournée est reçue de son équation HMT complète dans
\eqref{eq:mass-k-action} ; la vitesse est reçue des deux canaux
constitutifs de la section~\ref{sec:compat-lectores-formas}.
La relation circulaire convertit l’avancée relevée en rayon. Si
\(c=\widehat c\,u_C\), \(t_0=u_T\) et
\(u_L=u_Cu_T\), il en résulte
\(L_*/u_L=54\widehat c/(\pi_{\rm HMT}q_\alpha)\).
La référence radiale \(L_*\) et la base d’arête \(u_L\)
sont reliées avec leurs types dimensionnels explicites.

Dans une fibre finie, soit \(X=X^*>0\) un caractère positif inversible, \(U=2B\) et \(Y=UXU^*\). Le transport longitudinal est \(D_L=L_\alpha U\). Sa réponse radiale positive est définie au moyen de la métrique d'arrivée, \(\|R_Xv\|^2=\|XD_L^*v\|^2\) pour tout \(v\). L'énergie et la longueur conjuguée conservent le même caractère :
\[
 H_X=Q_{\rm clk}Y,\qquad M_X=H_X/c^2,\qquad
 \Lambda_X=L_\alpha Y^{-1}.
\]

\begin{theorem}[Coefficient radial commun et échelle aréale]
\label{rev5:rigidez-radial}
Les conditions précédentes déterminent une unique réponse positive
\(R_X=L_\alpha Y\). On a
\begin{equation}
 R_X\Lambda_X=L_\alpha^2I,\qquad
 H_X\Lambda_X=\hbar_{\rm ret}cI,\qquad
 R_X=\frac{L_\alpha^2}{\hbar_{\rm ret}c}H_X.
\end{equation}
Dans la réalisation physique qui identifie \(R_X\) au rayon
gravitationnel réduit \(R_X=(G/c^2)M_X\), le coefficient est
unique et commun à tous les caractères positifs admissibles :
\begin{equation}
 G=\frac{c^3L_\alpha^2}{\hbar_{\rm ret}},\qquad
 \frac{\hbar_{\rm ret}G}{c^3}=L_\alpha^2.
 \label{rev5:G-area}
\end{equation}
\end{theorem}
\begin{proof}
La polarisation de l'égalité des normes détermine \(R_X^2=D_LX^2D_L^*=L_\alpha^2UX^2U^*\). L'opérateur \(L_\alpha UXU^*\) est positif et son carré est celui indiqué. L'unicité de la racine positive prouve la première affirmation. Les deux identités de produits découlent d'une substitution. La convention physique donne \(L_\alpha Y=(GQ_{\rm clk}/c^4)Y\). L'inversibilité de \(Y\) permet de le simplifier, et l'expression de \(Q_{\rm clk}\) fournit \eqref{rev5:G-area}. Si un autre coefficient conserve les mêmes données, leur différence multipliée par \(M_X\) vaut zéro et est donc nulle.
\end{proof}

Le domaine de la preuve est la fibre positive finie ; il admet aussi
des opérateurs bornés uniformément positifs avec les mêmes identités.
La lecture du rayon gravitationnel réduit constitue la convention
physique déclarée de cette réalisation. La reconnaissance conventionnelle
ultérieure \(\ell_P^2=\hbar_{\rm ret}G/c^3\) identifie
\(\ell_P=L_\alpha\). La longueur s’obtient d’abord au moyen de
\eqref{rev5:longitud-generada} ; cette identification conserve son
constructeur. La formule circulaire équivalente est
\[
 G=\left(\frac{54}{\pi_{\rm HMT}}\right)^2
 \frac{c^5t_0^2}{\hbar_{\rm ret}}.
\]
Le facteur circulaire est déjà inclus lorsque l’on utilise \(L_\alpha\).
La longueur et la section d’action déterminent également
\[
 t_P=\frac{L_\alpha}{c},\qquad
 m_P=\frac{\hbar_{\rm ret}}{cL_\alpha},\qquad
 E_P=\frac{\hbar_{\rm ret}c}{L_\alpha}.
\]
Ces identités découlent de la substitution de \eqref{rev5:G-area} dans les
expressions conventionnelles des trois échelles ; la positivité
détermine leurs racines. Pour un mode de caractère \(y>0\), les
lectures sont \(m(y)=m_Py\), \(r_g(y)=L_\alpha y\) et
\(\bar\lambda(y)=L_\alpha/y\). Leur produit longitudinal
vaut \(r_g(y)\bar\lambda(y)=L_\alpha^2\). L’égalité des
deux longueurs équivaut à \(y=y^{-1}\) et, par positivité,
caractérise exactement \(y=1\). Le calendrier conserve
\(t_0=(\pi_{\rm HMT}/54)t_P\) et
\(108t_0=2\pi_{\rm HMT}t_P\).

Un rayon d’horizon \(R_h=L_\alpha\zeta_h\) conserve son
facteur propre \(\zeta_h=R_h/L_\alpha\). Son aire spatiale
appartient à l’horizon déclaré ; l’échelle élémentaire de l’opérateur
sectoriel suivant est \(L_\alpha^2\). Dans la spécialisation
de Schwarzschild pour la même masse, le rayon est \(2r_g(y)\)
et, par conséquent, \(\zeta_h=2y\) ; un horizon d’une autre réalisation
conserve sa condition géométrique de formation.
